\chapter{Bibliographical notes}
\label{chap:bibliographical-notes}

These notes are a reading map, not an additional review chapter.  The main
text was written to be calculationally self-contained; references here record
sources, alternative proofs, and routes to subjects which lie beyond the
book's connection-first line.  A work listed here is cited at the point where
it becomes useful, rather than being placed in an undifferentiated list of
``further references.''

\section*{Mathematical and quantum-mechanical foundations}

For Hilbert-space quantum mechanics, unbounded operators, and the spectral
theorem, the volumes of Reed and Simon remain a standard bridge between
physics and functional analysis \cite{ReedSimonI:1980}; Kato is the natural
next reference for closed operators, analytic families, and perturbation
theory \cite{Kato:1995}.  Trace ideals and determinant technology are treated
systematically by Simon \cite{SimonTraceIdeals:2005}.  Titchmarsh and Zettl
give complementary classical and modern routes through Sturm--Liouville
theory \cite{Titchmarsh:1962,Zettl:2005}.  The point of Stages I-A and I-B is
not to replace these books, but to derive exactly the subset needed when a
continuum spectral problem is analytically continued.

The direct-integral and operator-algebra discussion is indebted to the
decomposition theory initiated by \vN\ and to modern accounts of \vN\
algebras \cite{vonNeumann:1949Reduction,Takesaki:2002}.  Bratteli and Robinson
develop the $C^*$- and \vN-algebraic language in a form adapted to quantum
statistical mechanics \cite{BratteliRobinsonI:1987}; Haag is the standard
physical entrance to local quantum physics \cite{Haag:1996}.  These sources
also make clear why an energy-fiber direct integral, a central factor
decomposition, and a superselection decomposition answer different
questions.  The book keeps those three decompositions separate before using
them in scattering or gauge theory.

For rigged Hilbert spaces, the foundational distribution-theoretic setting
is due to Gelfand and collaborators \cite{Gelfand:1964}.  Pedagogical
treatments of continuous spectra and Gamow vectors may be found in
Refs.~\cite{DelaMadrid:2001wwa,DelaMadrid:2002cz}.  They are best read after
the explicit Gaussian continuation in Stage IV, where the topology on the
test space is seen to be part of the construction rather than decorative
notation.

\section*{Scattering, Jost functions, and long-range sectors}

Newton and Taylor are broad references for nonrelativistic scattering
\cite{Newton:1960,Taylor:1972}; Yafaev gives a mathematically systematic
account of wave operators and spectral methods \cite{Yafaev:1992}.  The
Zel'dovich school emphasizes resonances and quasistationary states
\cite{PerelomovZeldovich:1998}.  Rakityansky and Elander organize multichannel
scattering around the Jost matrix \cite{Rakityansky:2022Jost}, while the
Japanese textbook tradition represented here by Sasakawa provides a compact
route from elementary quantum mechanics to practical scattering
calculations \cite{Sasakawa:1991}.  Kukulin and collaborators connect this
material to resonant-state expansions and nuclear applications
\cite{Kukulin:1989}; the Berggren ensemble and its modern many-body use are
surveyed in Ref.~\cite{MichelPloszajczak:2021}.

Long-range scattering requires a different comparison dynamics, not a small
correction to a plane wave.  Dollard's construction is fundamental
\cite{Dollard:1964,Dollard:1971}.  For three charged particles, the
Faddeev--Merkuriev framework and subsequent asymptotic analyses display why
pairwise Coulomb phases do not by themselves give a single uniform global
form \cite{FaddeevMerkuriev:1993,AltMukhamedzhanov:1993,
MukhamedzhanovLieber:1996,Kadyrov:2003ThreeBody,Mukhamedzhanov:2006Leading,
Yakovlev:2022Coulomb,Keller:2009FaddeevMerkuriev}.  This is the few-body
counterpart of the infrared warning in field theory: the asymptotic state is
part of the dynamics.

The soft-photon problem begins historically with inclusive cancellation
\cite{BlochNordsieck:1937,YennieFrautschiSuura:1961}, but a state-level
description requires dressing.  Chung and Kulish--Faddeev are classic
constructions \cite{Chung:1965,KulishFaddeev:1970}; Buchholz explains the
charge-sector obstruction in algebraic terms \cite{Buchholz:1986}.  Modern
dressed-state and asymptotic-symmetry formulations relevant to the discussion
in Stage II include Refs.~\cite{ChenFrohlichPizzo:2010,
HiraiSugishita:2019,HiraiSugishita:2021,HiraiSugishita:2023}.  Dirac's
gauge-invariant charged field is an important precursor
\cite{Dirac:1955GaugeInvariant}.

For gauge constraints, BRST cohomology in the form used by particle
physicists is represented by Kugo and Ojima \cite{KugoOjima:1979}.  Gribov
copies and Singer's obstruction explain why the global nonperturbative
problem cannot be inferred from local gauge fixing alone
\cite{Gribov:1978,Singer:1978}.  Algebraic sector theory starts from
observables \cite{DoplicherHaagRoberts:1971,BuchholzFredenhagen:1982}.
Discussions of gauge-theory entanglement and edge structures make the loss of
a naive tensor product operational
\cite{CasiniHuertaRosabal:2014,DonnellyFreidel:2016,FewsterSplit:2015}; the
Hamiltonian lattice starting point goes back to Kogut and Susskind
\cite{KogutSusskind:1975}.  Quantum-simulation encodings give a current
computational realization of the same constraint issue \cite{PardoEtAl:2023}.

\section*{Exact WKB and resurgence}

Voros's analysis of the quartic oscillator is a landmark in the passage from
formal WKB to exact spectral information \cite{Voros:1983xx}.  The Japanese
school developed the microlocal and connection-formula foundations; useful
entry points include Refs.~\cite{AokiKawaiTakei1991RIMS853,Kawai:2005,
AokiIwakiTakahashi2019LoopType,AokiYoshida1993MicrolocalReduction}.  Problems
with many phases and global Stokes geometry are treated in
Refs.~\cite{AokiKawaiKoikeTakei2003GlobalAspects,
AokiKawaiKoikeTakei2004InfinitelyManyPhases,
AokiKawaiKoikeTakei2004MicroWKB}.  Hypergeometric benchmarks provide a useful
laboratory for checking normalization and connection conventions
\cite{AOKI:2021hmu}.

Delabaere and Pham give a classical resurgent treatment of semiclassical
asymptotics \cite{Delabaere:1999xx}; Dorigoni is a convenient field-theory
introduction to transseries and alien calculus \cite{Dorigoni:2014hea}.
Nikolaev's work gives a modern geometric view of exact WKB and
abelianization \cite{Nikolaev:2019img,Nikolaev:2021xzt,Nikolaev:2024}.
Wall-crossing formulations for quadratic differentials and related spectral
networks connect this material to wider geometric structures
\cite{Allegretti:2020dyt,Grassi:2022zuk}.  These developments are not required
to execute the one-dimensional connection calculations in Stage III, but
they clarify which parts of the construction survive for higher-order or
coupled systems.

The calculation program underlying the present book began by formulating
resonances nonperturbatively in exact WKB
\cite{MorikawaOgawa2025PTEP,Morikawa:2025vvs}, then unified Zel'dovich
regularization, complex scaling, and rigged Hilbert spaces
\cite{MorikawaOgawa2025JHEP,Morikawa:2025ezl}.  Continuum spectra and radial
equations were developed next
\cite{MorikawaOgawa2025Continuum,Morikawa:2025grx,Morikawa:2025inx}.  The
ordering of Stage III reflects the conceptual conclusion of that sequence:
the connection coefficient is primary; its operator and functional
realizations are representations of the same analytic datum.

\section*{Complex scaling, resonance functionals, and response}

The Aguilar--Balslev--Combes theorem originates in
Refs.~\cite{Aguilar:1971ve,Balslev:1971vb}; Simon's account isolates the
operator-theoretic mechanism \cite{Simon1972BalslevCombes}, and Herbst treats
important dilation-analytic spectral questions \cite{Herbst:1981ew}.
Exterior complex scaling is developed in Ref.~\cite{Simon1979ECS}.  Moiseyev's
book and reviews connect the theorem to practical non-Hermitian spectral
calculations \cite{Moiseyev:1998gjp,Moiseyev:2011,Moiseyev2011NHQM}.

Zel'dovich introduced Gaussian regularization for decaying states
\cite{Zel'dovich:1961}.  Berggren supplied the decisive resonant-state
normalization and completeness constructions
\cite{Berggren:1968zz,Berggren:1971,Berggren:1978,Berggren:1996}.  Stage IV
repeats the regulated integrations because the angular sector in which the
Gaussian limit exists is exactly what later becomes a complex-scaling or
exact-WKB Stokes-sector statement.

Feshbach's projection formalism is the source of the energy-dependent
effective Hamiltonian used throughout the book
\cite{Feshbach1958UnifiedI,Feshbach1962UnifiedII}.  Modern discussions of the
equivalence between effective Hamiltonians and resonant boundary conditions
include Ref.~\cite{Hatano2014Equivalence}.  Continuum level-density and
complex-scaled response methods in nuclear physics are reviewed in
Refs.~\cite{Myo:2014ypa,Myo:2023heu}.  Pseudospectral cautions should be read
alongside the elementary resolvent estimates; Trefethen's introduction is a
compact entry point \cite{Trefethen1999}.

\section*{Radial motion, few-body nuclei, and CDCC}

The Langer modification is traced to the original connection-formula paper
\cite{Langer:PhysRev.51.669}; the earlier geometric quantization context is
represented by Maslov and Arnol'd
\cite{Maslov1965PerturbationsAsymptotic,Arnold1967CharacteristicClassQuantization}.
Recent studies of power-law radial spectra provide useful numerical and
semiclassical comparisons \cite{delValle:2021pjv,delValle:2023feg}.  The
monodromy-renormalization proposal in Stage V is more specific: it keeps the
singular endpoint data invariant while changing the perturbative seed, and
its present limitations are stated in Sec.~\ref{subsec:open-monodromy-rg}.

CDCC grew from coupled-channel discretizations of breakup
\cite{Kamimura1986CDCC,Austern1987CDCC} and now has a broad reaction-theory
formulation \cite{Yahiro2012CDCC}.  Four-body CDCC and complex-scaled
Lippmann--Schwinger smoothing are developed in
Refs.~\cite{Matsumoto2004FourBodyCDCC,Kikuchi2009CSLS,Kikuchi2010CSLS}.
Gaussian expansion methods are described in Ref.~\cite{Hiyama:2003cu}; the
microscopic effective interaction used in many folding calculations is
represented by Ref.~\cite{Jeukenne1977JLM}.  The relation to exact few-body
equations is anchored by the Faddeev and Alt--Grassberger--Sandhas equations
\cite{Faddeev1961,AltGrassbergerSandhas1967}; the $R$-matrix interior--exterior
split is traced to Lane and Thomas \cite{LaneThomas1958RMatrix}.

The lithium and beryllium laboratories are not generic examples appended to
the formalism.  They record the line of nuclear-structure and reaction work
from which the model-space and channel questions arise
\cite{OgawaThesis:2022,arXiv:2003.05123,arXiv:2103.16865,
arXiv:2111.04285,Ogawa:2024He5,Ogawa:2025S18,Ogata:2026LiCDCC}.  Their role in
the book is to test whether pole positions, residues, smoothing, and measured
cross sections remain stable under one common truncation protocol.

\section*{Exceptional points}

Non-Hermitian effective Hamiltonians in open systems and their resonance
interpretation are reviewed by Rotter \cite{Rotter2007NonHermitian,
Rotter:2007zng}.  Broader modern perspectives are given in
Ref.~\cite{Ashida2020NonHermitian}.  Experiments and models which expose
encircling, state transfer, and topology include
Refs.~\cite{Naghiloo:2019dqw,Gong2018Topological,Gneiting:2020pew}.
Thresholds can modify exceptional-point decay laws and apparent order
\cite{Garmon2017NonExpThreshold,Garmon2021AnomalousEP}; this is why Stage VII
keeps the continuum regulator visible.

The complex-scaling calculation of resonance holonomy developed by the
authors treats an exceptional point as a multiple zero of the same incoming
connection coefficient used for an isolated resonance
\cite{Morikawa:2025xjq}.  Black-hole applications make the relation between
mode coalescence, excitation, and spectral instability especially current
\cite{PanossoMacedo:2025xnf,Takahashi:2025uwo}.  These papers motivate the
open problem of a regulator-independent continuum holonomy, but do not yet
settle it.

\section*{Black-hole quasinormal modes}

The Regge--Wheeler equation begins with Regge and Wheeler
\cite{Regge:1957td}; Zerilli supplied the even-parity formulation and charged
generalizations \cite{Zerilli:1970wzz,Zerilli:1974ai}.  Moncrief's
gauge-invariant Hamiltonian treatments
\cite{Moncrief:1974gw,Moncrief:1974ng,Moncrief:1974am} and the covariant
derivation of Martel and Poisson \cite{Martel:2005ir} are the main checks on
the derivation in Stage VIII.  Chandrasekhar's monograph remains the standard
source for transformations between master equations
\cite{Chandrasekhar:book,Chandrasekhar:1975nkd}.

General reviews of quasinormal modes include
Refs.~\cite{Kokkotas:1999bd,Berti:2009kk,Konoplya:2011qq}.  Leaver's continued
fractions and Green-function analysis are central numerical and analytic
benchmarks \cite{Leaver:1985ax,Leaver:1986gd}; Nollert developed the
high-overtone implementation \cite{Nollert:1993zz}.  Late-time tails and
branch-cut response are treated in Ref.~\cite{Ching:1994bd}.  Hyperboloidal
formulations clarify the operator and continuum structure
\cite{Ansorg:2016ztf,PanossoMacedo:2024nkw}, while pseudospectral instability
is emphasized in Ref.~\cite{Jaramillo:2020tuu}.

Exact-WKB approaches to black-hole spectra include
Refs.~\cite{Hatsuda:2019eoj,Hatsuda:2021gtn,Hatsuda:2023geo,HatsudaShiga:2026}.
The complex-scaling program developed in the final laboratories treats the
horizon conditions as rotated asymptotic sectors and retains continuum
response \cite{Ogawa:2026veu,Ogawa:2026dSCSM}.  The de Sitter extension makes
the two-horizon connection problem explicit.  For comparisons with current
spectroscopy and excitation questions, see
Refs.~\cite{Motohashi:2024fwt,DeAmicis:2025xuh,Berti:2025hly,Su:2026fvj}.

\section*{Resummation appendix}

The instanton and large-order background for
Chap.~\ref{chap:resummation-instantons} is presented pedagogically by
Mari\~no \cite{Marino:2015Instantons}.  Optimized perturbation, linear delta
expansion, variational perturbation, and ODM are represented by
Refs.~\cite{Stevenson:1981vj,Buckley:1992pc,Duncan:1992ba,
Kleinert:1995hc,Zinn-Justin:2010zzb}.  Multi-instanton structure in quantum
mechanics is developed in Refs.~\cite{Zinn-Justin:1981qzi,Jentschura:2011zza}.
The public calculation repository and Wolfram Community article give the
software provenance of the supplied compact code
\cite{Morikawa:ResummationRepository,Morikawa:2026WolframCommunity}.  The
book's conclusion is intentionally narrower than any one of these sources:
resummation becomes resonance theory only after its directional analytic
continuation is tied to a physical connection condition.
